\par\Needspace{4\baselineskip}
\begin{center}
\hypertarget{native-051-article033}{}\pdfbookmark[2]{33.}{article-033}33.
\end{center}
Thus, in conic sections whose eccentricity differs little from unity, that is, in ellipses and hyperbolas that closely approach a parabola, the methods described above for determining the true anomaly from the time and the time from the true anomaly\footnote{Since the time involves the factor $a^{\frac{3}{2}}$ or $b^{\frac{3}{2}}$, an error made in $M$ or $N$ is magnified all the more, the larger $a=\frac{p}{1-ee}$ or $b=\frac{p}{ee-1}$ is\raisebox{.45ex}{.}} do not allow all the precision one might wish for. Indeed, the inevitable errors grow as the orbit comes to resemble a parabola more closely and would ultimately exceed all bounds. Larger tables giving more than seven figures would reduce this uncertainty. It would nevertheless persist and exceed all bounds as soon as the orbit came too close to a parabola. Furthermore, the methods given above become rather laborious in this case, since some of them call for repeated indirect trials. This difficulty makes the work even more tedious if we use larger tables. It will therefore be worthwhile to develop a special method by which we can avoid this uncertainty and attain sufficient precision in this case with the aid of ordinary tables alone.

\par\Needspace{4\baselineskip}
\begin{center}
\hypertarget{native-051-article034}{}\pdfbookmark[2]{34.}{article-034}34.
\end{center}
The customary method for remedying these difficulties rests on the following principles. Consider an ellipse or hyperbola with eccentricity $e$, semiparameter $p$, and hence perihelion distance $=\frac{p}{1+e}=q$. \SourceLeaf{052}Let $t$ be the time after perihelion and $v$ the corresponding true anomaly. For the same time in a parabola with semiparameter $=2q$, or perihelion distance $=q$, let $w$ be the corresponding true anomaly, with the mass $\mu$ either neglected in both cases or assumed equal in both. We then clearly have
\[
\int\frac{pp\,\mathrm{d}v}{(1+e\cos v)^2}:\int\frac{4qq\,\mathrm{d}w}{(1+\cos w)^2}=\BareRoot{p}:\BareRoot{2q}
\]
with the integrals starting at $v=0$ and $w=0$, or
\[
\int\frac{(1+e)^{\frac{3}{2}}\mathrm{d}v}{(1+e\cos v)^2\BareRoot{2}}=\int\frac{2\mathrm{d}w}{(1+\cos w)^2}\raisebox{.45ex}{.}
\]
Denoting $\frac{1-e}{1+e}$ by $\alpha$ and $\operatorname{tang}\frac{1}{2}v$ by $\theta$, we obtain for the first integral
\[
\begin{aligned}
={}&\BareRoot{(1+\alpha)}\cdot\{\theta+\tfrac{1}{3}\theta^3(1-2\alpha)-\tfrac{1}{5}\theta^5(2\alpha-3\alpha\alpha)\\
&\qquad+\tfrac{1}{7}\theta^7(3\alpha\alpha-4\alpha^3)-\text{etc.}\},
\end{aligned}
\]
and for the second $=\operatorname{tang}\frac{1}{2}w+\frac{1}{3}\operatorname{tang}\frac{1}{2}w^3$. From this equation it is easy to determine $w$ from $\alpha$ and $v$, and $v$ from $\alpha$ and $w$, with the aid of infinite series. If preferred, we may replace $\alpha$ by $1-e=\frac{2\alpha}{1+\alpha}=\delta$. For $\alpha=0$ or $\delta=0$ we evidently have $v=w$, so these series will take the following form.
\[
\begin{aligned}
w&=v+\delta v^{\prime}+\delta\delta v^{\prime\prime}+\delta^3v^{\prime\prime\prime}+\text{etc.}\\
v&=w+\delta w^{\prime}+\delta\delta w^{\prime\prime}+\delta^3w^{\prime\prime\prime}+\text{etc.},
\end{aligned}
\]
Here $v^{\prime}$, $v^{\prime\prime}$, $v^{\prime\prime\prime}$, etc., are functions of $v$, and $w^{\prime}$, $w^{\prime\prime}$, $w^{\prime\prime\prime}$, etc., are functions of $w$. Whenever $\delta$ is a very small quantity, these series converge rapidly, and a few terms suffice to determine $w$ from $v$ or $v$ from $w$. From $w$ we find $t$, or $w$ from $t$, in the manner explained above for parabolic motion.
